\subsection{A fixed-field Frobenius bound for the remaining RH points}
\label{sec:rh-symplectic-spectrum}

The rank-four argument does not need the full source
\texttt{NumberSp4Alt} (with an integer simultaneously valid for
all admitted representations and all coefficient primes). The
original integral length comparison only needs a polynomial
bound for the sum of valuations at the actual coefficient prime,
on the actual family of irreducible points. We prove that
bound using a fixed finite collection of quadratic fields.
In particular the \texttt{WitA} witness is not needed for this
quantitative step.

\begin{theorem}[The actual rank-four trace estimate]
\label{thm:rh-symplectic-spectrum}
Retain the RH Case II datum of
Theorem~\ref{thm:rh-symplectic-arithmetic-data}, with coefficient
prime \(\ell\), host conductor \(N\), twist \(e\ge2\) and
auxiliary dyadic conductor \(2^n\). Set
\[
 M=\ell N2^n,\qquad k=(\ell-1)2^n.
\]
There is a choice of an imaginary quadratic field \(K\) from
the fixed list
\[
 \mathcal K=\{\mathbb Q(\sqrt{-7}),\mathbb Q(\sqrt{-11}),
             \mathbb Q(\sqrt{-19}),\mathbb Q(\sqrt{-43})\}
\]
and of a rational prime \(7\le p\le47\), split in \(K\), with
\(p\nmid M\). The field has a ramified prime outside the
ramification set of the actual rank-four datum, including \(\ell\).

Use one fixed coefficient isomorphism
\(\iota:\overline{\mathbb Q}_\ell\simeq\mathbb C\).
The only automorphy premise in this theorem is that each
irreducible actual point \(r\) under consideration has
\(r|_{G_K}\) attached, in the usual geometric algebraic
normalization, to a regular algebraic cuspidal representation
of \(\mathrm{GL}_4(\mathbb A_K)\). The source's
\texttt{IsAut4Alt} implies this premise. No separate
local compatibility witness is required.

Let \(\sigma\) be arithmetic Frobenius at the chosen \(p\).
For every finite family \(X\) of pairwise nonisomorphic
irreducible actual points satisfying that premise,
\[
 \begin{split}
  \operatorname{tr}r(\sigma)&\ne\operatorname{tr}r_0(\sigma)
                        &&(r\in X),\\
  (\log\ell)\sum_{r\in X}v_\ell\bigl(
        \operatorname{tr}r(\sigma)-\operatorname{tr}r_0(\sigma)\bigr)
       &\le C M^{130} k(e+4)^{13}\log p
        \le (\ell N(e+1))^{C'} .                         \tag{FS4}
 \end{split}
\]
Here \(v_\ell(\ell)=1\), and \(C,C'\) are absolute constants.
The choice of \(\sigma\) and the bound are independent of \(X\).
The same fixed-field automorphy premise suffices for
Theorem~\ref{thm:rh-symplectic-rigidity}.
\end{theorem}

\begin{proof}
\emph{A bounded choice of field and prime.}
The only odd primes dividing \(M\) belong to \(\{\ell,q,d\}\).
The four negative fundamental discriminants in \(\mathcal K\)
have distinct odd prime absolute values, so at least one of
them is prime to \(M\). The reduced primitive positive definite
forms of each of these discriminants consist of just one form:
\((1,1,2),(1,1,3),(1,1,5),(1,1,11)\), respectively. This is
checked by the reduction inequalities
\(|b|\le a\le c,\ 3a^2\le |\Delta_K|\), with the usual
nonnegative boundary convention. Thus each field has class
number one and unit group \(\{1,-1\}\).
The following four split primes in each field suffice:
\[
 \begin{array}{c|rrrr}
 \Delta_K&\multicolumn{4}{c}{\text{split primes}}\\ \hline
 -7&11&23&29&37\\
 -11&23&31&37&47\\
 -19&7&11&17&23\\
 -43&11&13&17&23 .
 \end{array}
\]
Their splitting follows from the quadratic residue symbol.
At most three of a row's primes divide \(M\).
Choose a remaining \(p\) and a place \(w\mid p\) in \(K\);
then \(K_w=\mathbb Q_p\).

Lemma~\ref{lem:rh-symplectic-cyclotomic-irreducibility} proves
absolute irreducibility on \(G_{K(\mu_{\ell^\infty})}\).
In particular all restrictions \(r|_{G_K}\) used here are
irreducible. The restriction map on their isomorphism classes
is injective: if two restrictions are isomorphic, their
one-dimensional intertwiner space is stable under \(G_{\mathbb Q}\),
with action factoring through \(\operatorname{Gal}(K/\mathbb Q)\).
Consequently the original representations differ by the trivial
or the nontrivial quadratic twist from this group. The latter
is impossible since both representations are unramified at
\(|\Delta_K|\), whereas that twist is ramified there. This
argument uses the actual ramification set, not a bound for
the degree of an auxiliary automorphy field.

\emph{A bounded algebraic character to remove the polarization twist.}
The similitude of \(r\) is
\(\tau=\varepsilon^{e+1}\delta\), where
\(\delta=\psi\xi\omega_T^{-1}\) has conductor dividing \(M\) and
\(\delta(-1)=(-1)^e\). Work with its complex character under
\(\iota\). The embedding
\((\mathbb Z/M\mathbb Z)^\times\hookrightarrow
(\mathcal O_K/M\mathcal O_K)^\times\) is injective.
A character on a subgroup of a finite abelian group extends
to the group (successively extend along cyclic quotients
using roots of unity). Choose an extension \(\vartheta\)
of \(\delta\) to this larger group. Its values have order
dividing the order of that group, which is at most \(M^2\).
Use a fixed embedding of \(K\) into \(\mathbb C\).
For a fractional ideal \(\mathfrak a=(a)\) prime to \(M\)
define the ideal character
\[
                 \nu(\mathfrak a)=a^{e+4}\vartheta(a\bmod M).
                                                               \tag{FS5}
\]
Here \(a\) is a unit at every prime dividing \(M\); reduction
is defined after clearing a denominator prime to \(M\).
This is independent of its generator because the only
units are \(\pm1\) and
\(\vartheta(-1)=(-1)^{e+4}\).
It is multiplicative and is the character on ideals of an
algebraic Hecke character of conductor dividing \(M\mathcal O_K\):
on a principal generator congruent to one modulo this ideal
it is exactly \(a^{e+4}\).
This also describes its infinity type. Its \(\ell\)-adic
realization \(\kappa\) is normalized by
\[
                 \iota\kappa(\operatorname{Frob}_v^{-1})
                        =\nu(v)\qquad(v\nmid \ell M).
\]
For the \(\ell\)-adic realization of an algebraic Hecke
character we use global reciprocity with the algebraic factors
at the embeddings above \(\ell\) removed from the finite
idele character. In particular its local Weil--Deligne
inertia at every place is the finite local character of
\(\nu\) or its inverse according to the reciprocal convention.
Its conductor divides \(M\mathcal O_K\). At \(\ell\) this
statement follows by separating the algebraic local factors
from the finite character: the former are crystalline
characters at each embedding and do not contribute to
Weil--Deligne inertia. This is the rank-one local comparison
used for the algebraic characters in the earlier auxiliary
tower constructions.

For a prime \(v\nmid \ell M\), the identity
\(\nu(v)\nu(\bar v)=(Nv)^{e+4}\delta(Nv)\)
follows directly from (FS5).
Since \(\varepsilon(\operatorname{Frob}_v^{-1})=(Nv)^{-1}\)
and \(\iota\delta(\operatorname{Frob}_v^{-1})=\delta(Nv)^{-1}\)
(with the right side the finite Dirichlet character), Chebotarev
gives on \(G_K\)
\[
                  \kappa\kappa^c=\tau^{-1}\varepsilon^{-3}.
                                                               \tag{FS6}
\]
In particular the point \(r'=r|_{G_K}\otimes\kappa\)
has conjugate duality \(r'{}^c\simeq r'{}^\vee\varepsilon^{-3}\).
The arithmetic Frobenius convention in the final trace
comparison is unchanged.

The assumed automorphic attachment of \(r|_{G_K}\) can be
twisted by the algebraic Hecke character realizing \(\kappa\).
This preserves cuspidality and regular algebraicity, since
at each embedding it shifts all four labeled weights by
the same integer. Its resulting representation \(\Pi\)
is conjugate self-dual: (FS6) gives the required local
contragredient identities at almost all places in the
geometric normalization \(\operatorname{rec}(\Pi_v
\otimes|\det|^{(1-4)/2})\).
Strong multiplicity one gives \(\Pi^\vee\simeq\Pi^c\).
By Chebotarev the attached semisimple Galois representation
is exactly \(r'\). This twisting argument uses the
cuspidal attachment in the premise, not an assumed
\texttt{WitA} or a new automorphy lifting theorem.

\emph{The conductor and weight of the counting family.}
The actual local conditions in
Proposition~\ref{prop:rh-symplectic-local-finite} give
potential crystallinity at \(\ell\), and its four graded
inertia characters are those of
\((\tau,\Psi,\alpha,1)\).
Their finite parts have conductor dividing \(M\);
away from \(\ell\) the actual residual inertia kernel
condition likewise makes the representation factor through
the finite cyclotomic inertia quotient with modulus \(M\).
For the latter assertion use the equality of the
prime-to-\(\ell\) finite character kernels before and
after reduction. The base cyclotomic powers are
unramified away from \(\ell\).
Restriction to \(K\) preserves these conductor bounds:
local reciprocity and the local norm send
\(1+M\mathcal O_{K_v}\) into \(1+M\mathbb Z_s\) at a
place \(v\mid s\).
Twisting by \(\kappa\) multiplies each finite inertia
character by another character trivial on
\(1+M\mathcal O_{K_v}\).
Its algebraic part at \(\ell\) is potentially crystalline,
so the monodromy of \(r'\) is zero at all finite places.
The exact conductor bound for its Weil--Deligne representation
is therefore
\[
                  a_v(r')\le4v(M).
\]
Full local--global compatibility for conjugate self-dual
cohomological cuspidal representations, including the
monodromy at the coefficient prime, is available by
Caraiani's Theorem~1.1 in each of the two papers
\cite{support-spectrum-Caraiani2012,support-spectrum-Caraiani2014}.
In particular the argument does not require the
Shin-regular weight condition; the weights at \(e=2\)
need not be Shin-regular.
The conductor exponent of \(\Pi_v\) obeys the same bound.
By the local new-vector theorem and its correction
\cite{support-spectrum-newvector}, \(\Pi\) has a nonzero
fixed vector under the principal congruence subgroup
\[
 K_4(M^4)=\ker\bigl(\mathrm{GL}_4(\widehat{\mathcal O}_K)
         \longrightarrow\mathrm{GL}_4(\mathcal O_K/M^4)\bigr).
                                                               \tag{FS7}
\]
This includes the place above the coefficient prime; using
only semisimplified local compatibility there would not
justify (FS7).

The labeled Hodge--Tate multisets of \(r|_{G_K}\) are the
four distinct integers \(-(e+1),-e,-1,0\) at both embeddings,
using the convention \(\operatorname{HT}(\varepsilon)=-1\).
The labeled weights of \(r'\) are obtained by adding the
rank-one weight of \(\kappa\), with the two shifts \(e+4\)
and zero. Equivalently their multisets, up to interchanging
the embeddings, are
\[
           \{e+4,e+3,4,3\},\qquad \{0,-1,-e,-e-1\}.
                                                               \tag{FS8}
\]
Thus they lie in \([-e-4,e+4]\), and all pairwise
differences within either embedding are at most \(e+1\).
The relation between the Hodge--Tate labels and the
cohomological highest weights is the one in
BLGGT Theorem~2.1.1(3) \cite{support-spectrum-BLGGT}.
In particular the two possible highest weight types
are determined by (FS8) and permutation of embeddings.
Their coefficient representation has dimension at
most \((e+4)^{12}\) by the Weyl dimension formula
(six positive roots at each embedding). Its action on
the positive split real center is trivial, by conjugate
self-duality: the sums of the two highest weight types
cancel. This permits the standard locally symmetric
space with its real split center divided out.

Let \(\mathcal A\) be the entire finite indexed family of
conjugate self-dual regular algebraic cuspidal representations
of \(\mathrm{GL}_4/K\) of either indicated weight type with
a fixed vector under (FS7). The cardinality counts
representations, not just their individual trace values.
Fix a neat finite level \(U_0\) once for each of the four
fields, and set \(U=U_0\cap K_4(M^4)\).
The Borel--Serre cell complex for the fixed level lifts
to that at \(U\) with covering degree at most
\[
        [U_0:U]\le |\mathrm{GL}_4(\mathcal O_K/M^4)|
                                      \le M^{128}.
\]
The sum of the cellular cochain dimensions is at most
this degree times the fixed number of cells times the
dimension of the coefficient representation.
Cuspidal cohomology injects into ordinary cohomology,
and every member of \(\mathcal A\) contributes a nonzero
summand, distinguished by its Hecke eigensystem.
Using the Borel--Serre and Clozel results already applied
in Theorem~\ref{thm:support-spectrum}, we obtain
\[
                     |\mathcal A|\le C_K M^{128}(e+4)^{12}.
                                                               \tag{FS9}
\]
Taking the maximum over the four fields makes \(C_K\)
absolute. Clozel's rationality and Galois conjugation
theorem preserves the cohomological level and permutes
the labeled weights; it acts on the Hecke eigenvalues
in their algebraic normalization. Therefore the family
of attached geometric Frobenius multisets at \(w\),
indexed by \(\mathcal A\), is algebraic and
Galois-stable, with multiplicities. The injectivity of
restriction proved above and the invertibility of
the single twist \(\kappa\) ensure that distinct original
points give distinct indices in this family.

\emph{A nonzero integer controlling all actual trace differences.}
For \(\Pi\in\mathcal A\) let
\(\lambda_{\Pi,1},\ldots,\lambda_{\Pi,4}\) be the
geometrically normalized Frobenius eigenvalues at \(w\).
They all have absolute value \(p^{3/2}\) under every
complex embedding, by the conjugate self-dual temperedness
theorem (Caraiani 2012, Theorem~1.2) and the geometric
normalization. The companion Galois representations
at every rational coefficient prime supplied by
BLGGT Theorem~2.1.1 show that these eigenvalues and
their inverses are integral away from \(p\):
each is an eigenvalue of an invertible matrix in a
stable integral lattice at a good prime.
At \(p\), the companion representation is crystalline
since \(\Pi_w\) is unramified, by part~(4) of that
same theorem. Its Hodge--Tate weights lie in
\([-e-4,e+4]\) by (FS8). Weak admissibility bounds all
its Newton slopes in this same interval. To see the
bound, extend the coefficient field to split Frobenius;
on a rank-one Frobenius eigenline the induced Hodge
number is at least the minimum Hodge number, so its
Newton number is at least that minimum. Apply the same
argument to the dual for the upper bound.
Here \(K_w=\mathbb Q_p\), so the bound concerns the
valuation normalized by \(v_p(p)=1\).
The same argument applies for every coefficient embedding
above \(p\); the geometric Frobenius and the crystalline
Frobenius normalizations can differ by inversion, which
does not affect this symmetric bound.
Consequently \(p^{e+4}\lambda_{\Pi,i}^{-1}\) is an
algebraic integer.

Let \(\mathcal V\) be the Galois orbit of the single
algebraic integer \(\nu(w)\).
Every member has modulus \(p^{(e+4)/2}\), and
\(|\mathcal V|\le2M^2\), by the definition (FS5).
For \(\Pi\in\mathcal A\) and \(z\in\mathcal V\) set
\[
                     t_{\Pi,z}=z\sum_i\lambda_{\Pi,i}^{-1}.
\]
These are algebraic, with \(p^{e+4}t_{\Pi,z}\) integral,
and
\[
                       |t_{\Pi,z}|\le4p^{(e+1)/2}.
                                                               \tag{FS10}
\]
For an actual representation the value
\(\iota\operatorname{tr}r(\sigma)\) is exactly
\(t_{\Pi,\nu(w)}\), not the geometric trace:
\(\sigma\) restricts to an arithmetic Frobenius at \(w\)
because \(p\) splits in \(K\), and the inverse of the
geometric Frobenius acts with eigenvalues
\(\lambda_{\Pi,i}^{-1}\). Untwisting multiplies this
arithmetic trace by \(\nu(w)\).

Let \(\mathcal E\) denote the Galois orbit of the triple
\[
       (z_1,z_2,z_3)=
       (\iota\delta(p),\iota\psi(p),
                         \iota(\xi\omega_T^{-1})(p)).
\]
Here all three are evaluated as Dirichlet character
values at the integer \(p\). Their orders divide \(k\)
since \(2^n\ge N\) and the odd host character orders
divide two. Thus \(|\mathcal E|\le k\).
For a triple \(z'\in\mathcal E\), put
\[
                   T(z')=p^{e+1}z'_1+p^ez'_2+pz'_3+1 .
\]
For every such triple, \(|T(z')|\ge
p^{e+1}-p^e-p-1>(p-2)p^e>4p^{(e+1)/2}\), since
\(p\ge7\) and \(e\ge2\). This proves nonvanishing
of every difference with (FS10).
Form the finite product with its indicated multiplicities,
\[
 B=\prod_{\Pi\in\mathcal A}\ \prod_{z\in\mathcal V}\
      \prod_{z'\in\mathcal E}
            p^{e+4}\bigl(t_{\Pi,z}-T(z')\bigr).
                                                               \tag{FS11}
\]
The empty product is one.
It is a nonzero algebraic integer fixed by every
Galois automorphism of \(\overline{\mathbb Q}/\mathbb Q\),
hence a nonzero rational integer. The absolute value
of each factor is at most \(p^{4(e+4)}\), by (FS10) and
the triangle inequality. By (FS9),
\[
        \log|B|\le 8C_K M^{130} k(e+4)^{13}\log p.
\]
For any actual family \(X\), the factors associated to
its distinct indices, the fixed \(z=\nu(w)\), and the
original triple, are exactly the actual trace differences
times \(p^{e+4}\). Every other factor has nonnegative
\(\ell\)-valuation by integrality. Since \(\ell\ne p\),
the total valuation of these actual trace differences
is bounded by \(v_\ell(B)\le\log|B|/\log\ell\).
This proves the first bound in (FS4). The second follows
from (BU7), \(M\le128 N^2\ell^6\),
\(k\le128N\ell^6\), and the fixed finite set of primes.

Finally the proof of Theorem~\ref{thm:rh-symplectic-rigidity}
uses the automorphic attachment only over a field whose
restriction satisfies (SR2). The present \(K\) has that
property. The explicit twist (FS6) also gives its
conjugate self-dual normalization without a
\texttt{WitA} witness; twisting does not alter the
adjoint. Thus the same proof of local finite conditions
and Thorne's theorem proves rigidity under precisely
the present fixed-field premise.
\end{proof}

The theorem replaces the quantitative rank-four automorphic
spectrum supplier for the RH specialization by (FS4). It does
not prove the existence of the cuspidal attachments in its
premise, and does not assert the entire simultaneous
\texttt{NumberSp4Alt} proposition for arbitrary admitted data.
For the original RH Case II integral comparison it supplies
the required trace bound once the remaining point arithmetic
and graded length interfaces are established.
